\documentclass[11pt,a4paper]{article}
% Spolecny styl pro ucitelske pripravy MAT (Work/MAT)
% Verze 1.1, 2026-10-23
\NeedsTeXFormat{LaTeX2e}
\ProvidesPackage{priprava}[2026/10/23 v1.1 Sablona ucitelskych priprav MAT]

% Zakladni typografie a cestina
\RequirePackage[T1]{fontenc}
\RequirePackage[utf8]{inputenc}
\RequirePackage[czech]{babel}
\RequirePackage{lmodern}
\RequirePackage{microtype}
\RequirePackage{amsmath,amssymb}
\RequirePackage{xcolor}
\RequirePackage{enumitem}
\RequirePackage[most]{tcolorbox}
\RequirePackage{tabularx,booktabs,array}
\RequirePackage{fancyhdr}
\RequirePackage{needspace}
\RequirePackage{geometry}
\RequirePackage[hidelinks]{hyperref}

\geometry{a4paper,top=17mm,bottom=18mm,left=23mm,right=23mm}
\definecolor{notegray}{gray}{0.39}
\setlength{\parindent}{0pt}
\setlength{\parskip}{2.5pt plus 1pt minus 0.5pt}
\setlist[itemize]{topsep=2pt,itemsep=1pt,parsep=0pt,leftmargin=1.45em}
\setlist[enumerate]{topsep=2pt,itemsep=2pt,parsep=0pt,leftmargin=1.65em}
\setcounter{secnumdepth}{0}
\renewcommand{\arraystretch}{1.12}

% Zahlavi: vlevo rocnik | globalni tema, vpravo cislo: nazev hodiny
% Zapati zustava zcela prazdne.
\pagestyle{fancy}
\fancyhf{}
\renewcommand{\headrulewidth}{0.35pt}
\renewcommand{\footrulewidth}{0pt}
\setlength{\headheight}{14pt}
\newcommand{\nastavPripravu}[4]{%
  \fancyhead[L]{\small #1\enspace\textbar\enspace #2}%
  \fancyhead[R]{\small #3: #4}%
}

% Minimalisticke cernobile boxy. Kazda definice zvlast.
\tcbset{commonlesson/.style={enhanced,breakable=false,colback=white,colframe=black,
  colbacktitle=white,coltitle=black,boxrule=0.45pt,
  arc=0.8mm,outer arc=0.8mm,left=3mm,right=3mm,top=1.6mm,bottom=1.6mm,
  toptitle=1mm,bottomtitle=0.7mm,fonttitle=\bfseries,
  attach boxed title to top left={xshift=2.5mm,yshift=-1.15mm},
  boxed title style={colback=white,colframe=white,boxrule=0pt,left=1mm,right=1mm}}}
\newtcolorbox{zadaniBox}[1][]{commonlesson,title={Zadání},#1}
\newtcolorbox{prikladBox}[1][]{commonlesson,title={Příklad},#1}
\newtcolorbox{definiceBox}[1][]{commonlesson,title={Definice},#1}

% Spolecny box pro zadani a jeho reseni, bez mezery mezi castmi.
% Pouziti:
%   \begin{prikladReseniBox}
%     Zadani ulohy.
%     \tcblower
%     \begin{align*} ... \end{align*}
%   \end{prikladReseniBox}
% Spodni cast ma jemny sedy podklad a delici linku; lze v ni sazet
% vice radku matematiky. Obe casti maji jen jeden spolecny ramecek.
\tcbset{solutionlesson/.style={
  commonlesson,
  enhanced jigsaw,
  breakable,
  colbacklower=black!3,
  segmentation style={solid,draw=black!25,line width=0.35pt},
  fontlower=\small\color{notegray},
  before lower={\noindent\textbf{Řešení:}\enspace},
  leftlower=3mm,rightlower=3mm
}}
\newtcolorbox{prikladReseniBox}[1][]{solutionlesson,title={Příklad},#1}
\newtcolorbox{zadaniReseniBox}[1][]{solutionlesson,title={Zadání},#1}

% Pomucky pro autora pripravy: sedive poznamky; \reseni pro starsi pripravy
\newcommand{\poznamka}[1]{\par\begingroup\small\color{notegray}\itshape
  \textnormal{\textbf{Poznámka pro učitele.}}\enspace #1\par\endgroup}
\newcommand{\reseni}[1]{\par\begingroup\small\color{notegray}
  \textbf{Očekávané řešení:}\enspace #1\par\endgroup}
\newcommand{\blok}[1]{\needspace{5\baselineskip}\subsection*{#1}}
\newcommand{\vyrok}[1]{\uv{#1}}

\nastavPripravu{1. ročník}{Algebraické výrazy}{42}{Sčítání a odčítání}
\begin{document}
\poznamka{45 minut: 5 min opakování členů; 8 min výklad a první vzor; 8 min rozdíl; 12 min samostatná práce; 8 min společná kontrola; 4 min rezerva. U více proměnných kontroluj shodné exponenty všech písmen.}
\blok{Které členy můžeme sečíst?}
Ze zkušenosti se známými jednotkami víme, že $3\,\mathrm m+2\,\mathrm m=5\,\mathrm m$, ale $3\,\mathrm m+2\,\mathrm{m}^2$ nelze sloučit v jedinou veličinu. Obdobně $3x^2+2x^2=5x^2$, zatímco $3x^2+2x$ neslučujeme.
\begin{definiceBox}[title={Definice: podobné členy}]
\textbf{Podobné (stejnorodé) členy} mají stejnou proměnnou část, včetně všech exponentů. Sčítáme a odčítáme pouze jejich \textbf{koeficienty}: $ax^ky^\ell+bx^ky^\ell=(a+b)x^ky^\ell$.
\end{definiceBox}
\begin{definiceBox}[title={Definice: opačný mnohočlen}]
Opačný mnohočlen k $P$ je $-P=(-1)P$. Mění se znaménko \textbf{každého} členu a platí $P+(-P)=0$.
\end{definiceBox}
\poznamka{Zdůvodni $-(A+B)=(-1)(A+B)=-A-B$ distributivitou. Plus před závorkou členy nemění; mínus násobí celou závorku $-1$.}
\blok{Řešené příklady}
\begin{prikladReseniBox}[title={Příklad 1: součet tří mnohočlenů}]
Sečti $A=3x^2+1$, $B=2x^2+xy-3$ a $C=x+y$.
\tcblower
$A+B+C=3x^2+1+2x^2+xy-3+x+y=\boxed{5x^2+xy+x+y-2}$.
\end{prikladReseniBox}
\begin{prikladReseniBox}[title={Příklad 2: rozdíl a změna znamének}]
Uprav $(4x^3-2x^2+x-5)-(x^3+3x^2-2x-7)$.
\tcblower
$4x^3-2x^2+x-5-x^3-3x^2+2x+7=\boxed{3x^3-5x^2+3x+2}$.
\end{prikladReseniBox}
\begin{prikladReseniBox}[title={Příklad 3: proč se součet opačných rovná nule?}]
Je-li $F(x)=4x^4+7x^3-2x-2$, uprav $F(x)+[-F(x)]$.
\tcblower
$F+(-F)=4x^4+7x^3-2x-2-4x^4-7x^3+2x+2=\boxed0$.
\end{prikladReseniBox}
\blok{Samostatné procvičování}
\begin{zadaniReseniBox}[title={Samostatná práce 1: základní úpravy}]
\begin{enumerate}[label=\alph*)]
\item $(2x^2-3x+1)+(x^2+x-4)$.
\item $(4a^2-2a+5)-(a^2+3a-1)$.
\item $3xy-2x^2+5xy+x^2$.
\item $(3xy-2x^2+y^2)-(xy+x^2-4y^2)$.
\end{enumerate}
\tcblower
\begin{enumerate}[label=\alph*)]
\item $3x^2-2x-3$.
\item $3a^2-5a+6$.
\item $8xy-x^2$.
\item $2xy-3x^2+5y^2$.
\end{enumerate}
\end{zadaniReseniBox}
\begin{zadaniReseniBox}[title={Samostatná práce 2: vnořené závorky}]
\begin{enumerate}[label=\alph*)]
\item $4x^2-[2x^2-(3x-1)]+(x^2-2x)$.
\item $(x^3-2x+1)-[(x^3+x-4)-(2x-3)]$.
\item $2x-[3x-(x-4)]$.
\end{enumerate}
\tcblower
\begin{enumerate}[label=\alph*)]
\item $3x^2+x-1$.
\item $-x+2$.
\item $-4$.
\end{enumerate}
\end{zadaniReseniBox}
\poznamka{Nech žáky ověřit výsledky dosazením například $x=0$ a $x=1$. Zkouška umí odhalit chybu, není však sama o sobě důkazem identity.}
\blok{Rezerva a domácí práce}
\begin{zadaniReseniBox}[title={Rozšiřující příklad 1: zlomkové koeficienty}]
Uprav $\left(\frac12a^2-\frac34a+1\right)+\left(\frac13a^2+\frac14a-2\right)-\frac16a^2$.
\tcblower
$\left(\frac12+\frac13-\frac16\right)a^2+\left(-\frac34+\frac14\right)a-1=\boxed{\frac23a^2-\frac12a-1}$.
\end{zadaniReseniBox}
\begin{zadaniReseniBox}[title={Domácí úkol 1}]
\begin{enumerate}[label=\alph*)]
\item $(3x^2-2x+4)-(x^2+5x-1)$.
\item $-(2a^2-a+3)+(a^2-4a+2)$.
\item $2ab-3a^2-(5ab-a^2)$.
\item $3x-[2x-(x+5)]$.
\item $(2x^2-3xy+y^2)-(x^2+xy-2y^2)$.
\end{enumerate}
\tcblower
\begin{enumerate}[label=\alph*)]
\item $2x^2-7x+5$.
\item $-a^2-3a-1$.
\item $-3ab-2a^2$.
\item $2x+5$.
\item $x^2-4xy+3y^2$.
\end{enumerate}
\end{zadaniReseniBox}
\end{document}
